% GENERATED FILE. Edit canonical lessons, then run npm run generate:books.
% canonical-source-hash: fnv1a64:2d38594f799b63b2
% canonical-lessons: ML-C390-yatrayayappu, ML-C390-vitavannuka, ML-C390-yatrayakkuka, ML-C390-vitaparayuka, ML-C390-natuvituka

\chapter{Farewells}
\label{ch:ml-farewells}
\hlchaptermodality{390}

\begin{chapteropening}
\textbf{By the end of this chapter:} \emph{I can say a send-off, to bid farewell, to see off, to say goodbye and to leave one's homeland.}

Complete the last lesson of chapter 390: I can say a send-off, to bid farewell, to see off, to say goodbye and to leave one's homeland.
\end{chapteropening}

\section[\texorpdfstring{\ml{യാത്രയയപ്പ്} (yātrayayappŭ)}{യാത്രയയപ്പ് (yātrayayappŭ)}]{\ml{യാത്രയയപ്പ്} (yātrayayappŭ) — a send-off; a farewell party}

\label{lesson:ML-C390-yatrayayappu}



\begin{quote}
Before the new one: say the Malayalam for to settle, then the Malayalam for to migrate.
\end{quote}

\subsection*{You'll want to know: \ml{യാത്രയയപ്പ്}}
\textbf{\ml{യാത്രയയപ്പ്}} — \emph{yātrayayappŭ} — \textquotedblleft{}a send-off; a farewell party\textquotedblright{}.

\begin{quote}
\ml{ഓഫീസിൽ} \ml{മാനേജർക്ക്} \ml{ഒരു} \ml{യാത്രയയപ്പ്} \ml{ഉണ്ടായിരുന്നു}.
\end{quote}

\begin{quote}
\emph{ōphīsil mānējarkkŭ oru yātrayayappŭ uṇṭāyirunnu.}
\end{quote}

\begin{quote}
There was a send-off for the manager at the office.
\end{quote}

\ml{യാത്ര}, a journey, + \ml{അയപ്പ്}, a sending.

\begin{cousinweb}
\textbf{\ml{യാത്രയയപ്പ്}} (\emph{yātrayayappŭ}) joins \ml{യാത്ര} (\emph{yātra}), \textquotedblleft{}a journey\textquotedblright{}, borrowed from Sanskrit \emph{yātrā} (CDIAL 10456), and the native \ml{അയപ്പ്}, \textquotedblleft{}a sending\textquotedblright{}, the noun of \ml{അയയ്ക്കുക}, \textquotedblleft{}to send\textquotedblright{}.
\end{cousinweb}

\begin{grammarlens}[title={What you've built}]
One more word for telling what happened.
\end{grammarlens}

\subsection*{Your turn}
\begin{itemize}
\raggedright
  \item \emph{Say it:} \emph{yātrayayappŭ}
  \item \emph{Say it:} \emph{yātrayayappŭ}, once more
  \item \emph{Say it:} \emph{ōphīsil mānējarkkŭ oru yātrayayappŭ uṇṭāyirunnu.} — There was a send-off for the manager at the office.
  \item \emph{Recall:} say the Malayalam for to settle, then the Malayalam for to migrate, then say \emph{yātrayayappŭ} again
\end{itemize}

\begin{tcolorbox}[breakable,colback=teal!4,colframe=teal!35!black,title={Before you move on}]
What does \ml{യാത്രയയപ്പ്} mean? (\textquotedblleft{}A send-off; a farewell party\textquotedblright{}.) Say it once more.
\end{tcolorbox}
\section[\texorpdfstring{\ml{വിടവാങ്ങുക} (viṭavāṅṅuka)}{വിടവാങ്ങുക (viṭavāṅṅuka)}]{\ml{വിടവാങ്ങുക} (viṭavāṅṅuka) — to bid farewell; to take leave}

\label{lesson:ML-C390-vitavannuka}



\begin{quote}
Before the new one: say the Malayalam for to migrate, then the Malayalam for a send-off.
\end{quote}

\subsection*{You'll want to know: \ml{വിടവാങ്ങുക}}
\textbf{\ml{വിടവാങ്ങുക}} — \emph{viṭavāṅṅuka} — \textquotedblleft{}to bid farewell; to take leave\textquotedblright{}.

\begin{quote}
\ml{മുപ്പത്} \ml{വർഷത്തിന്} \ml{ശേഷം} \ml{അധ്യാപകൻ} \ml{പള്ളിക്കൂടത്തിൽ} \ml{നിന്ന്} \ml{വിടവാങ്ങി}.
\end{quote}

\begin{quote}
\emph{muppatŭ varṣattinŭ śēṣaṁ adhyāpakan paḷḷikkūṭattil ninnŭ viṭavāṅṅi.}
\end{quote}

\begin{quote}
After thirty years the teacher took leave of the school.
\end{quote}

Past: \ml{വിടവാങ്ങി}, \emph{viṭavāṅṅi}.

\ml{വിട}, leave, + \ml{വാങ്ങുക}, to take: to take one's leave. Also a polite word for to die.

\begin{cousinweb}
\textbf{\ml{വിടവാങ്ങുക}} (\emph{viṭavāṅṅuka}) is a native Dravidian compound of \ml{വിട} (\emph{viṭa}), \textquotedblleft{}leave, a farewell\textquotedblright{}, from \ml{വിടുക}, \textquotedblleft{}to let go\textquotedblright{} (DEDR 5393), and \ml{വാങ്ങുക} (\emph{vāṅṅuka}), \textquotedblleft{}to take, to receive\textquotedblright{} (DEDR 5336).
\end{cousinweb}

\begin{grammarlens}[title={What you've built}]
One more word for telling what happened.
\end{grammarlens}

\subsection*{Your turn}
\begin{itemize}
\raggedright
  \item \emph{Say it:} \emph{viṭavāṅṅuka}
  \item \emph{Say it:} \emph{viṭavāṅṅuka}, once more
  \item \emph{Say it:} \emph{muppatŭ varṣattinŭ śēṣaṁ adhyāpakan paḷḷikkūṭattil ninnŭ viṭavāṅṅi.} — After thirty years the teacher took leave of the school.
  \item \emph{Recall:} say the Malayalam for to migrate, then the Malayalam for a send-off, then say \emph{viṭavāṅṅuka} again
\end{itemize}

\begin{tcolorbox}[breakable,colback=teal!4,colframe=teal!35!black,title={Before you move on}]
What does \ml{വിടവാങ്ങുക} mean? (\textquotedblleft{}To bid farewell; to take leave\textquotedblright{}.) Say it once more.
\end{tcolorbox}
\section[\texorpdfstring{\ml{യാത്രയാക്കുക} (yātrayākkuka)}{യാത്രയാക്കുക (yātrayākkuka)}]{\ml{യാത്രയാക്കുക} (yātrayākkuka) — to see off}

\label{lesson:ML-C390-yatrayakkuka}



\begin{quote}
Before the new one: say the Malayalam for a send-off, then the Malayalam for to bid farewell.
\end{quote}

\subsection*{You'll want to know: \ml{യാത്രയാക്കുക}}
\textbf{\ml{യാത്രയാക്കുക}} — \emph{yātrayākkuka} — \textquotedblleft{}to see off\textquotedblright{}.

\begin{quote}
\ml{ഞങ്ങൾ} \ml{അമ്മാവനെ} \ml{സ്റ്റേഷനിൽ} \ml{യാത്രയാക്കി}.
\end{quote}

\begin{quote}
\emph{ñaṅṅaḷ ammāvane sṟṟēṣanil yātrayākki.}
\end{quote}

\begin{quote}
We saw uncle off at the station.
\end{quote}

Past: \ml{യാത്രയാക്കി}, \emph{yātrayākki}.

The traveller takes -\ml{എ}: \ml{അവനെ} \ml{യാത്രയാക്കി}.

\begin{cousinweb}
\textbf{\ml{യാത്രയാക്കുക}} (\emph{yātrayākkuka}) joins \ml{യാത്ര} (\emph{yātra}), \textquotedblleft{}a journey\textquotedblright{}, borrowed from Sanskrit \emph{yātrā} (CDIAL 10456), and the native \ml{ആക്കുക}, \textquotedblleft{}to make\textquotedblright{}: to set someone on their journey.
\end{cousinweb}

\begin{grammarlens}[title={What you've built}]
One more word for telling what happened.
\end{grammarlens}

\subsection*{Your turn}
\begin{itemize}
\raggedright
  \item \emph{Say it:} \emph{yātrayākkuka}
  \item \emph{Say it:} \emph{yātrayākkuka}, once more
  \item \emph{Say it:} \emph{ñaṅṅaḷ ammāvane sṟṟēṣanil yātrayākki.} — We saw uncle off at the station.
  \item \emph{Recall:} say the Malayalam for a send-off, then the Malayalam for to bid farewell, then say \emph{yātrayākkuka} again
\end{itemize}

\begin{tcolorbox}[breakable,colback=teal!4,colframe=teal!35!black,title={Before you move on}]
What does \ml{യാത്രയാക്കുക} mean? (\textquotedblleft{}To see off\textquotedblright{}.) Say it once more.
\end{tcolorbox}
\section[\texorpdfstring{\ml{വിടപറയുക} (viṭapaṟayuka)}{വിടപറയുക (viṭapaṟayuka)}]{\ml{വിടപറയുക} (viṭapaṟayuka) — to say goodbye}

\label{lesson:ML-C390-vitaparayuka}



\begin{quote}
Before the new one: say the Malayalam for to bid farewell, then the Malayalam for to see off.
\end{quote}

\subsection*{You'll want to know: \ml{വിടപറയുക}}
\textbf{\ml{വിടപറയുക}} — \emph{viṭapaṟayuka} — \textquotedblleft{}to say goodbye\textquotedblright{}.

\begin{quote}
\ml{കുട്ടികൾ} \ml{മുത്തശ്ശിയോട്} \ml{വിടപറഞ്ഞു}.
\end{quote}

\begin{quote}
\emph{kuṭṭikaḷ muttaśśiyōṭŭ viṭapaṟaññu.}
\end{quote}

\begin{quote}
The children said goodbye to grandmother.
\end{quote}

Past: \ml{വിടപറഞ്ഞു}, \emph{viṭapaṟaññu}.

The person takes -\ml{ഓട്}. In the news, also to pass away.

\begin{cousinweb}
\textbf{\ml{വിടപറയുക}} (\emph{viṭapaṟayuka}) is a native Dravidian compound of \ml{വിട} (\emph{viṭa}), \textquotedblleft{}leave, a farewell\textquotedblright{} (DEDR 5393), and \ml{പറയുക} (\emph{paṟayuka}), \textquotedblleft{}to say\textquotedblright{} (DEDR 4031).
\end{cousinweb}

\begin{grammarlens}[title={What you've built}]
One more word for telling what happened.
\end{grammarlens}

\subsection*{Your turn}
\begin{itemize}
\raggedright
  \item \emph{Say it:} \emph{viṭapaṟayuka}
  \item \emph{Say it:} \emph{viṭapaṟayuka}, once more
  \item \emph{Say it:} \emph{kuṭṭikaḷ muttaśśiyōṭŭ viṭapaṟaññu.} — The children said goodbye to grandmother.
  \item \emph{Recall:} say the Malayalam for to bid farewell, then the Malayalam for to see off, then say \emph{viṭapaṟayuka} again
\end{itemize}

\begin{tcolorbox}[breakable,colback=teal!4,colframe=teal!35!black,title={Before you move on}]
What does \ml{വിടപറയുക} mean? (\textquotedblleft{}To say goodbye\textquotedblright{}.) Say it once more.
\end{tcolorbox}
\section[\texorpdfstring{\ml{നാടുവിടുക} (nāṭuviṭuka)}{നാടുവിടുക (nāṭuviṭuka)}]{\ml{നാടുവിടുക} (nāṭuviṭuka) — to leave one's homeland}

\label{lesson:ML-C390-natuvituka}



\begin{quote}
Before the new one: say the Malayalam for to see off, then the Malayalam for to say goodbye.
\end{quote}

\subsection*{You'll want to know: \ml{നാടുവിടുക}}
\textbf{\ml{നാടുവിടുക}} — \emph{nāṭuviṭuka} — \textquotedblleft{}to leave one's homeland\textquotedblright{}.

\begin{quote}
\ml{ജോലി} \ml{തേടി} \ml{അവൻ} \ml{ചെറുപ്പത്തിൽ} \ml{നാടുവിട്ടു}.
\end{quote}

\begin{quote}
\emph{jōli tēṭi avan ceṟuppattil nāṭuviṭṭu.}
\end{quote}

\begin{quote}
He left his homeland young, in search of work.
\end{quote}

Past: \ml{നാടുവിട്ടു}, \emph{nāṭuviṭṭu}.

\ml{നാട്}, homeland, + \ml{വിടുക}, to leave.

\begin{cousinweb}
\textbf{\ml{നാടുവിടുക}} (\emph{nāṭuviṭuka}) is a native Dravidian compound of \ml{നാട്} (\emph{nāṭŭ}), \textquotedblleft{}the country, one's homeland\textquotedblright{} (DEDR 3638), and \ml{വിടുക} (\emph{viṭuka}), \textquotedblleft{}to leave, to quit\textquotedblright{} (DEDR 5393).
\end{cousinweb}

\begin{grammarlens}[title={What you've built}]
One more word for telling what happened.
\end{grammarlens}

\subsection*{Your turn}
\begin{itemize}
\raggedright
  \item \emph{Say it:} \emph{nāṭuviṭuka}
  \item \emph{Say it:} \emph{nāṭuviṭuka}, once more
  \item \emph{Say it:} \emph{jōli tēṭi avan ceṟuppattil nāṭuviṭṭu.} — He left his homeland young, in search of work.
  \item \emph{Recall:} say the Malayalam for to see off, then the Malayalam for to say goodbye, then say \emph{nāṭuviṭuka} again
\end{itemize}

\begin{tcolorbox}[breakable,colback=teal!4,colframe=teal!35!black,title={Before you move on}]
What does \ml{നാടുവിടുക} mean? (\textquotedblleft{}To leave one's homeland\textquotedblright{}.) Say it once more.
\end{tcolorbox}
